\chapter{1977年湖南卷}

\section{解答题}

\begin{problem}
\begin{enumerate}
    \item \( \log_x 3\frac{3}{8}=-1\frac{1}{2} \)，求 \( x \).
    \item \(\lg 4+\lg 9+2\sqrt{\left(\lg 6\right)^2-2\lg 6+1}=?\)
    \item 求函数 \( y=\frac{\sqrt{3-x}}{\sqrt{x-2}} \) 的定义域（即自变量 \( x \) 的取值范围）.
    \item 解不等式：\(x^2-3x-28<0\).
    \item 证明 \(\cos 15^\circ+\sin 15^\circ=\frac{\sqrt{6}}{2}\).
    \item 求中心在原点，焦点在 \(x\) 轴上，且经过 \(A\left(5,0\right)\)，\(B\left(4,\frac{12}{5}\right)\) 两点的椭圆方程.
\end{enumerate}
\end{problem}

\begin{solution}
\begin{enumerate}
\item 对数的底数必须满足 \(x>0\) 且 \(x\ne1\)，并且
\(3\frac{3}{8}=\frac{27}{8}\)，\(-1\frac{1}{2}=-\frac{3}{2}\). 由对数定义，得
\(x^{-\frac{3}{2}}=\frac{27}{8}=\left(\frac{3}{2}\right)^3\). 两边取倒数并开立方，得 \(x^{\frac{3}{2}}=\left(\frac{2}{3}\right)^3\)，所以 \(x^{\frac{1}{2}}=\frac{2}{3}\)，从而 \(x=\frac{4}{9}\). 该值满足 \(x>0\) 且 \(x\ne1\)，故原方程的解为 \(x=\frac{4}{9}\).

\item 因为 \(6<10\)，所以 \(\lg 6<1\). 因而
\(\sqrt{\left(\lg 6\right)^2-2\lg 6+1}=\sqrt{\left(\lg 6-1\right)^2}=1-\lg 6\). 于是原式
\(=\lg 4+\lg 9+2-2\lg 6=\lg 36+2-\lg 36=2\).

\item 分式有意义，必须使分子根式的被开方数非负，同时使分母根式的被开方数严格为正，因此 \(x\) 满足
\[
\begin{cases}
3-x\ge0,\\
x-2>0.
\end{cases}
\]
由第一式得 \(x\le3\)，由第二式得 \(x>2\)，两者取交集得 \(2<x\le3\). 所以函数的定义域为 \((2,3]\).

\item 将多项式因式分解，得
\(x^2-3x-28=(x-7)(x+4)\). 因此原不等式等价于 \((x-7)(x+4)<0\). 两个因式的零点为 \(-4\) 和 \(7\)，在 \(x<-4\)、\(-4<x<7\)、\(x>7\) 三个区间上，乘积的符号依次为正、负、正；又因不等式是严格不等式，两个零点不能取. 所以解集为 \((-4,7)\)，即 \(-4<x<7\).

\item 设 \(u=\cos15^\circ+\sin15^\circ\). 则
\[
u^2=\cos^2 15^\circ+\sin^2 15^\circ+2\sin15^\circ\cos15^\circ
=1+\sin30^\circ=1+\frac12=\frac32
\]
由于 \(15^\circ\) 在第一象限，\(\cos15^\circ>0\) 且 \(\sin15^\circ>0\)，故 \(u>0\). 因而 \(u=\sqrt{\frac32}=\frac{\sqrt6}{2}\)，即
\(\cos15^\circ+\sin15^\circ=\frac{\sqrt6}{2}\).

也可以利用和差化积公式：
\[
\cos15^\circ+\sin15^\circ=\sin75^\circ+\sin15^\circ
=2\sin45^\circ\cos30^\circ
\]
\[
2\sin45^\circ\cos30^\circ
=2\cdot\frac{\sqrt2}{2}\cdot\frac{\sqrt3}{2}
=\frac{\sqrt6}{2}
\]

\item 因为椭圆的中心在原点且焦点在 \(x\) 轴上，可设其标准方程为
\(\frac{x^2}{a^2}+\frac{y^2}{b^2}=1\)，其中 \(a>b>0\). 点 \(A(5,0)\) 在椭圆上，所以
\(\frac{25}{a^2}=1\)，从而 \(a=5\). 将点 \(B\left(4,\frac{12}{5}\right)\) 代入，得
\(\frac{4^2}{25}+\frac{\left(\frac{12}{5}\right)^2}{b^2}=1\)，\(\frac{16}{25}+\frac{144}{25b^2}=1\).
因此 \(\frac{144}{25b^2}=\frac9{25}\)，即 \(b^2=16\). 由 \(b>0\)，得 \(b=4\). 所以所求椭圆方程为
\(\frac{x^2}{25}+\frac{y^2}{16}=1\).
\end{enumerate}
\end{solution}

\begin{problem}
某县农机厂金工车间共有 \(86\) 个工人，已知每个工人平均每天可以加工甲种部件 \(15\) 个，或者乙种部件 \(12\) 个，或者丙种部件 \(9\) 个. 问应安排加工甲种部件，乙种部件和丙种部件各多少人，才能使加工后的 \(3\) 个甲种部件，\(2\) 个乙种部件和 \(1\) 个丙种部件恰好配套（即加工后的甲种部件个数是丙种部件个数的 \(3\) 倍，乙种部件个数是丙种部件个数的 \(2\) 倍）.
\end{problem}

\begin{solution}
设安排加工甲、乙、丙三种部件的工人数分别为 \(x\)、\(y\)、\(z\). 工人总数为 \(86\)，而甲、乙、丙三种部件的产量分别为 \(15x\)、\(12y\)、\(9z\)，所以由配套条件得
\[
\begin{cases}
x+y+z=86,\\
15x=3\cdot9z,\\
12y=2\cdot9z.
\end{cases}
\]
将后两式分别约去公因数，化为
\[
\begin{cases}
x+y+z=86,\\
5x-9z=0,\\
2y-3z=0.
\end{cases}
\]
将第一式乘以 \(2\) 后减去第三式，得 \(2x+5z=172\). 再将这个等式乘以 \(5\)，减去第二式的 \(2\) 倍，得
\(43z=860\)，所以 \(z=20\). 代入 \(2y-3z=0\)，得 \(y=30\). 再由第一式得 \(x=86-30-20=36\). 三个数均为非负整数，且代回原条件可得产量分别为 \(540\)、\(360\)、\(180\)，恰为 \(3:2:1\). 因此应安排甲、乙、丙三种部件工人分别为 \(36\) 人、\(30\) 人、\(20\) 人.
\end{solution}

\begin{problem}
已知方程 \(ax^2-(2a+1)x+(a+1)=0\).
\begin{enumerate}
    \item 有一个根为 \(0\)；
    \item 两个根的和为 \(0\)；
    \item 两个根的乘积为 \(2\).
\end{enumerate}
分别求 \(a\) 的值.
\end{problem}

\begin{solution}
由于题目给出的是二次方程，必须有 \(a\ne0\).
\begin{enumerate}
\item 若 \(0\) 是方程的一个根，将 \(x=0\) 代入得 \(a+1=0\)，所以 \(a=-1\). 此时二次项系数不为零，条件确实成立.

\item 设方程的两个根为 \(x_1\)、\(x_2\). 由韦达定理，
\(x_1+x_2=\frac{2a+1}{a}\). 当两根的和为 \(0\) 时，\(\frac{2a+1}{a}=0\). 因为 \(a\ne0\)，所以 \(2a+1=0\)，得 \(a=-\frac12\). 此时方程为 \(-\frac12x^2+\frac12=0\)，确有两个根 \(1\)、\(-1\)，条件成立.

\item 同样由韦达定理，\(x_1x_2=\frac{a+1}{a}\). 令其等于 \(2\)，得 \(a+1=2a\)，所以 \(a=1\). 此时方程为 \(x^2-3x+2=0=(x-1)(x-2)\)，两个根的乘积确为 \(2\).
\end{enumerate}
\end{solution}

\begin{problem}
把表面积分别是 \(36\pi\mathrm{~cm}^2\)、\(64\pi\mathrm{~cm}^2\) 和 \(100\pi\mathrm{~cm}^2\) 的三个锡球熔成一个大锡球，求这个大锡球的半径.
\end{problem}

\begin{solution}
设大锡球的半径为 \(R\)，三个小锡球的半径分别为 \(r_1\)、\(r_2\)、\(r_3\). 由球的表面积公式 \(S=4\pi r^2\)，有
\(r_1=\sqrt{\frac{36\pi}{4\pi}}=3\)，\(r_2=\sqrt{\frac{64\pi}{4\pi}}=4\)，\(r_3=\sqrt{\frac{100\pi}{4\pi}}=5\).
熔化过程中没有材料损失，因此体积守恒. 由球的体积公式 \(V=\frac43\pi r^3\)，得
\[
\frac43\pi R^3=\frac43\pi\left(r_1^3+r_2^3+r_3^3\right)
\]
约去 \(\frac43\pi\)，并代入三个半径，得
\(R^3=3^3+4^3+5^3=27+64+125=216\). 由于半径为正，\(R=6\). 所以这个大锡球的半径为 \(6\mathrm{~cm}\).
\end{solution}

\begin{problem}
已知四边形 \(ABCD\) 中，\(\angle A\)、\(\angle B\)、\(\angle C\)、\(\angle D\) 成等差数列，且 \(\angle C=3\angle A\)，求 \(\angle A\)、\(\angle B\)、\(\angle C\)、\(\angle D\) 的度数.
\end{problem}

\begin{solution}
设等差数列的首项为 \(x\)，公差为 \(d\)，则
\(\angle A=x\)、\(\angle B=x+d\)、\(\angle C=x+2d\)、\(\angle D=x+3d\). 由 \(\angle C=3\angle A\)，得 \(x+2d=3x\)，即 \(d=x\). 所以四个角依次为 \(x\)、\(2x\)、\(3x\)、\(4x\). 四边形内角和为 \(360^\circ\)，故
\(x+2x+3x+4x=360^\circ\)，解得 \(x=36^\circ\). 因而
\(\angle A=36^\circ\)、\(\angle B=72^\circ\)、\(\angle C=108^\circ\)、\(\angle D=144^\circ\).
\end{solution}

\begin{problem}
在 \(\triangle ABC\) 中，用 \(a\)、\(b\)、\(c\) 依次表示 \(\angle A\)、\(\angle B\)、\(\angle C\) 的对边，已知 \(\angle B=45^\circ\)、\(b=10\)、\(c=5\sqrt6\)，不查表求 \(C\)、\(A\)、\(a\).
\end{problem}

\begin{solution}
由正弦定理，\(\frac{c}{\sin C}=\frac{b}{\sin B}\)，所以
\[
\sin C=\frac{c\sin B}{b}
=\frac{5\sqrt6\cdot\frac{\sqrt2}{2}}{10}
=\frac{\sqrt3}{2}
\]
三角形内角满足 \(0^\circ<C<180^\circ\)，故 \(\sin C=\frac{\sqrt3}{2}\) 给出两种可能：\(C=60^\circ\) 或 \(C=120^\circ\).

当 \(C=60^\circ\) 时，\(A=180^\circ-45^\circ-60^\circ=75^\circ\). 再由正弦定理，
\[
a=\frac{b\sin A}{\sin B}
=\frac{10\sin75^\circ}{\sin45^\circ}
=\frac{10\cdot\frac{\sqrt6+\sqrt2}{4}}{\frac{\sqrt2}{2}}
=5(\sqrt3+1)
\]

当 \(C=120^\circ\) 时，\(A=180^\circ-45^\circ-120^\circ=15^\circ\). 同理
\[
a=\frac{b\sin A}{\sin B}
=\frac{10\sin15^\circ}{\sin45^\circ}
=\frac{10\cdot\frac{\sqrt6-\sqrt2}{4}}{\frac{\sqrt2}{2}}
=5(\sqrt3-1)
\]
两种情形下的第三角均为正，均能构成三角形. 因此答案为
\(C=60^\circ\)，\(A=75^\circ\)，\(a=5(\sqrt3+1)\)，或
\(C=120^\circ\)，\(A=15^\circ\)，\(a=5(\sqrt3-1)\).
\end{solution}

\begin{problem}
已知圆内接四边形 \(ABCD\) 中，\(\angle A=60^\circ\)、\(\angle B=90^\circ\)、\(AB=2\)、\(CD=1\). 求 \(BC\) 和 \(DA\).

\bitmapfigure[max width=4.8cm]{img/1977/hunan/q7_fig1.png}
\end{problem}

\begin{solution}
方法一：由圆周角定理，\(\angle B=90^\circ\) 所对的弦 \(AC\) 是外接圆的直径. 因此同弦 \(AC\) 所对的圆周角 \(\angle D=90^\circ\). 四边形内角和为 \(360^\circ\)，所以 \(\angle C=360^\circ-\angle A-\angle B-\angle D=120^\circ\).

设 \(BC=x>0\)、\(DA=y>0\). 在直角三角形 \(ABC\) 和 \(ADC\) 中，斜边都是 \(AC\)，由勾股定理分别得 \(AC^2=AB^2+BC^2=4+x^2\) 与 \(AC^2=AD^2+DC^2=y^2+1\)，从而
\(4+x^2=y^2+1\).

对三角形 \(ABD\) 使用余弦定理，得 \(BD^2=AB^2+AD^2-2AB\cdot AD\cos60^\circ=4+y^2-2y\). 对三角形 \(BCD\) 使用余弦定理，得 \(BD^2=BC^2+CD^2-2BC\cdot CD\cos120^\circ=x^2+1+x\). 比较两个 \(BD^2\) 的表达式，得 \(4+y^2-2y=x^2+1+x\).

由第一式得 \(x^2=y^2-3\)，代入第二式，得 \(4+y^2-2y=y^2-2+x\)，即 \(x+2y=6\). 因而 \(x=6-2y\). 代入 \(x^2=y^2-3\)，得
\[
(6-2y)^2=y^2-3
\quad\Longleftrightarrow\quad
3y^2-24y+39=0
\quad\Longleftrightarrow\quad
(y-4)^2=3
\]
所以 \(y=4\pm\sqrt3\). 当 \(y=4+\sqrt3\) 时，\(x=6-2y=-2-2\sqrt3<0\)，不符合边长为正的条件；故取 \(y=4-\sqrt3\)，此时 \(x=6-2y=2\sqrt3-2\). 因此
\(BC=2\sqrt3-2\)，\(DA=4-\sqrt3\).

方法二：由上面的圆周角关系，\(\angle D=90^\circ\)、\(\angle C=120^\circ\). 过 \(B\) 作 \(BE\perp AD\)，垂足为 \(E\)；过 \(C\) 作 \(CF\perp BE\)，垂足为 \(F\). 因为 \(BE\perp AD\) 且 \(CF\perp BE\)，而 \(CD\perp AD\)，所以四边形 \(CDEF\) 为矩形，进而 \(ED=CF\)、\(FE=CD=1\).

\solutionfigure{\bitmapfigure[max width=4.9cm]{img/1977/hunan/q7_solution_fig1.png}}

在直角三角形 \(ABE\) 中，\(\angle ABE=90^\circ-\angle A=30^\circ\)，故
\(AE=AB\sin30^\circ=2\cdot\frac12=1\)，且 \(BE=AB\cos30^\circ=2\cdot\frac{\sqrt3}{2}=\sqrt3\).

在直角三角形 \(BCF\) 中，\(BF=BE-FE=\sqrt3-1\). 又 \(\angle CBF=90^\circ-\angle ABE=60^\circ\)，所以
\(BC=\frac{BF}{\cos60^\circ}=2(\sqrt3-1)\). 另外
\(CF=BF\tan60^\circ=(\sqrt3-1)\sqrt3=3-\sqrt3\). 因为 \(AD=AE+ED=AE+CF\)，所以 \(AD=1+3-\sqrt3=4-\sqrt3\). 这与方法一所得结果一致.
\end{solution}

\section{参考题}

\begin{problem}
用数学归纳法或其他方法证明：\(1^4+2^4+3^4+\dots+n^4=\frac{n(n+1)(2n+1)(3n^2+3n-1)}{30}\).
\end{problem}

\begin{solution}
方法一：数学归纳法. 设命题 \(P(n)\) 为
\(1^4+2^4+\dots+n^4=\frac{n(n+1)(2n+1)(3n^2+3n-1)}{30}\)，其中 \(n\) 为正整数.

当 \(n=1\) 时，左边为 \(1^4=1\)，右边为
\(\frac{1\cdot2\cdot3\cdot5}{30}=1\)，所以 \(P(1)\) 成立.

假设 \(P(k)\) 成立，即
\(1^4+2^4+\dots+k^4=\frac{k(k+1)(2k+1)(3k^2+3k-1)}{30}\). 则
\[
\begin{aligned}
1^4+2^4+\dots+k^4+(k+1)^4
&=\frac{k(k+1)(2k+1)(3k^2+3k-1)+30(k+1)^4}{30}\\
&=\frac{(k+1)(6k^4+39k^3+91k^2+89k+30)}{30}\\
&=\frac{(k+1)(k+2)(6k^3+27k^2+37k+15)}{30}\\
&=\frac{(k+1)(k+2)(2k+3)(3k^2+9k+5)}{30}\\
&=\frac{(k+1)((k+1)+1)(2(k+1)+1)\left(3(k+1)^2+3(k+1)-1\right)}{30}.
\end{aligned}
\]
这里用到的因式分解可直接验证：
\(6k^4+39k^3+91k^2+89k+30=(k+2)(2k+3)(3k^2+9k+5)\). 因此 \(P(k+1)\) 成立. 由数学归纳法，对任意正整数 \(n\)，原等式成立.

方法二：利用递推式. 对正整数 \(r\) 记 \(S_r=1^r+2^r+\dots+n^r\). 对任意正整数 \(k\)，由二项式定理，对每个 \(j=1\)，\(2\)，\(\dots\)，\(n\) 有
\[
(j+1)^{k+1}-j^{k+1}-1
=\mathrm C_{k+1}^1j^k+\mathrm C_{k+1}^2j^{k-1}+\dots+\mathrm C_{k+1}^kj
\]
对 \(j=1\) 到 \(n\) 求和，左边望远镜相消，得到
\[
\mathrm C_{k+1}^1S_k+\mathrm C_{k+1}^2S_{k-1}+\dots+\mathrm C_{k+1}^kS_1
=(n+1)^{k+1}-(n+1)
\]
依次令 \(k=1\)，\(2\)，\(3\)，\(4\)，可得
\(S_1=\frac{n(n+1)}2\)，\(S_2=\frac{n(n+1)(2n+1)}6\)，\(S_3=\frac{n^2(n+1)^2}{4}\).
当 \(k=4\) 时，代入前三式，得
\[
\begin{aligned}
S_4
&=\frac{(n+1)^5-(n+1)-5S_1-10S_2-10S_3}{5}\\
&=\frac{(n+1)^5-(n+1)-\frac52n(n+1)-\frac53n(n+1)(2n+1)-\frac52n^2(n+1)^2}{5}\\
&=\frac{n(n+1)(2n+1)(3n^2+3n-1)}{30}.
\end{aligned}
\]
所以 \(S_4=1^4+2^4+\dots+n^4\) 等于题中所给表达式.
\end{solution}

\begin{problem}
计算下列积分：
\begin{enumerate}
\item \(\int_0^1 xe^{-x}\,dx\)；
\item \(\int_0^a x^2\sqrt{a^2-x^2}\,dx\)（\(a>0\)）.
\end{enumerate}
\end{problem}

\begin{solution}
\begin{enumerate}
\item 由分部积分，取 \(u=x\)，\(dv=e^{-x}\,dx\)，则 \(du=dx\)，\(v=-e^{-x}\). 因而
\[
\int_0^1 xe^{-x}\,dx
=\left[-xe^{-x}\right]_0^1+\int_0^1e^{-x}\,dx
=-\frac1{\e}+\left[-\e^{-x}\right]_0^1
=1-\frac2{\e}
\]

\item 因为 \(a>0\)，令 \(x=a\sin t\)，其中 \(0\le t\le\frac\pi2\). 则 \(dx=a\cos t\,dt\)，且
\(\sqrt{a^2-x^2}=\sqrt{a^2-a^2\sin^2t}=a\cos t\). 当 \(x=0\) 时 \(t=0\)，当 \(x=a\) 时 \(t=\frac\pi2\)，所以
\[
\begin{aligned}
\int_0^a x^2\sqrt{a^2-x^2}\,dx
&=a^4\int_0^{\frac\pi2}\sin^2t\cos^2t\,dt
=\frac{a^4}{4}\int_0^{\frac\pi2}\sin^2 2t\,dt\\
&=\frac{a^4}{8}\int_0^{\frac\pi2}(1-\cos4t)\,dt
=\frac{a^4}{8}\left[t-\frac14\sin4t\right]_0^{\frac\pi2}
=\frac{a^4\pi}{16}.
\end{aligned}
\]
\end{enumerate}
\end{solution}
